% ============================================================
%  MTH636 -- Game Theory | Lecture 29
%  Strategic Form to Extensive Form & Incomplete Information Games
% ============================================================
\documentclass[12pt, a4paper]{article}

% ============================================================
%  MTH636 -- Game Theory
%  Notes Preamble — Color-Coded, TikZ-Enhanced
% ============================================================

% --- Core packages ---
\usepackage[T1]{fontenc}
\usepackage[utf8]{inputenc}
\usepackage{lmodern}
\usepackage[margin=2.5cm, top=2.8cm, bottom=2.8cm]{geometry}
\usepackage{microtype}

% --- Math ---
\usepackage{amsmath, amssymb, amsthm, mathtools}
\usepackage{bm}        % bold math
\usepackage{cancel}    % strike-through in math

% --- Color & Boxes ---
\usepackage[dvipsnames, svgnames, x11names, table]{xcolor}
\usepackage[most]{tcolorbox}
\tcbuselibrary{breakable, skins, theorems}

% --- TikZ & Diagrams ---
\usepackage{tikz}
\usepackage{tikz-qtree}
\usetikzlibrary{arrows.meta, positioning, shapes.geometric, fit, backgrounds,
                decorations.pathreplacing, calc, matrix, trees}
\usepackage{pgfplots}
\pgfplotsset{compat=1.18}

% --- Tables ---
\usepackage{booktabs}
\usepackage{multirow}
\usepackage{array}
\usepackage{tabularx}
\usepackage{makecell}

% --- Typography & Layout ---
\usepackage{enumitem}
\usepackage{parskip}
\usepackage{fancyhdr}
\usepackage{titlesec}
\usepackage{titletoc}
\usepackage{hyperref}
\usepackage{graphicx}
\usepackage{float}
\usepackage{caption}

% --- Color Palette ---
\definecolor{defblue}{HTML}{1A5276}      % Definition header
\definecolor{defbluebg}{HTML}{D6EAF8}    % Definition background
\definecolor{thmgreen}{HTML}{1E8449}     % Theorem header
\definecolor{thmbg}{HTML}{D5F5E3}        % Theorem background
\definecolor{exorange}{HTML}{B7770D}     % Example header
\definecolor{exbg}{HTML}{FEF9E7}         % Example background
\definecolor{warnred}{HTML}{922B21}      % Warning/Key Insight header
\definecolor{warnbg}{HTML}{FADBD8}       % Warning background
\definecolor{noticegray}{HTML}{2C3E50}   % Note header
\definecolor{noticebg}{HTML}{EBF5FB}     % Note background
\definecolor{taketawayyellow}{HTML}{7D6608} % Key Takeaway header
\definecolor{takeawaybg}{HTML}{FDFEFE}   % Takeaway background
\definecolor{accentpurple}{HTML}{6C3483} % Accent / proof
\definecolor{lightgray}{HTML}{F2F3F4}    % shading

% ============================================================
%  TCOLORBOX ENVIRONMENTS
% ============================================================

% -- Definition (Blue) --
\newtcolorbox{defbox}[2][]{
  enhanced, breakable,
  colback=defbluebg, colframe=defblue, coltitle=white,
  fonttitle=\bfseries\sffamily,
  title={Definition: #2},
  attach boxed title to top left={yshift=-2mm, xshift=4mm},
  boxed title style={colback=defblue, rounded corners},
  arc=4pt, boxrule=1.2pt,
  left=6pt, right=6pt, top=8pt, bottom=6pt,
  #1
}

% -- Theorem / Proposition (Green) --
\newtcolorbox{thmbox}[2][]{
  enhanced, breakable,
  colback=thmbg, colframe=thmgreen, coltitle=white,
  fonttitle=\bfseries\sffamily,
  title={Theorem: #2},
  attach boxed title to top left={yshift=-2mm, xshift=4mm},
  boxed title style={colback=thmgreen, rounded corners},
  arc=4pt, boxrule=1.2pt,
  left=6pt, right=6pt, top=8pt, bottom=6pt,
  #1
}

% -- Proposition (Green variant) --
\newtcolorbox{propbox}[2][]{
  enhanced, breakable,
  colback=thmbg, colframe=thmgreen!70!black, coltitle=white,
  fonttitle=\bfseries\sffamily,
  title={Proposition: #2},
  attach boxed title to top left={yshift=-2mm, xshift=4mm},
  boxed title style={colback=thmgreen!70!black, rounded corners},
  arc=4pt, boxrule=1.2pt,
  left=6pt, right=6pt, top=8pt, bottom=6pt,
  #1
}

% -- Example (Orange/Yellow) --
\newtcolorbox{exbox}[2][]{
  enhanced, breakable,
  colback=exbg, colframe=exorange, coltitle=white,
  fonttitle=\bfseries\sffamily,
  title={Example: #2},
  attach boxed title to top left={yshift=-2mm, xshift=4mm},
  boxed title style={colback=exorange, rounded corners},
  arc=4pt, boxrule=1.2pt,
  left=6pt, right=6pt, top=8pt, bottom=6pt,
  #1
}

% -- Key Insight / Warning (Red) --
\newtcolorbox{warnbox}[2][]{
  enhanced, breakable,
  colback=warnbg, colframe=warnred, coltitle=white,
  fonttitle=\bfseries\sffamily,
  title={Key Insight: #2},
  attach boxed title to top left={yshift=-2mm, xshift=4mm},
  boxed title style={colback=warnred, rounded corners},
  arc=4pt, boxrule=1.2pt,
  left=6pt, right=6pt, top=8pt, bottom=6pt,
  #1
}

% -- Note (Gray-Blue) --
\newtcolorbox{notebox}[1][]{
  enhanced, breakable,
  colback=noticebg, colframe=noticegray!60, coltitle=white,
  fonttitle=\bfseries\sffamily,
  title={Note},
  attach boxed title to top left={yshift=-2mm, xshift=4mm},
  boxed title style={colback=noticegray!60, rounded corners},
  arc=4pt, boxrule=0.8pt,
  left=6pt, right=6pt, top=8pt, bottom=6pt,
  #1
}

% -- Key Takeaway --
\newtcolorbox{takebox}[1][]{
  enhanced, breakable,
  colback=takeawaybg, colframe=taketawayyellow, coltitle=white,
  fonttitle=\bfseries\sffamily,
  title={Key Takeaway},
  attach boxed title to top left={yshift=-2mm, xshift=4mm},
  boxed title style={colback=taketawayyellow, rounded corners},
  arc=4pt, boxrule=1.2pt,
  left=6pt, right=6pt, top=8pt, bottom=6pt,
  #1
}

% -- Intuition --
\newtcolorbox{intubox}[1][]{
  enhanced, breakable,
  colback=accentpurple!8, colframe=accentpurple!60, coltitle=white,
  fonttitle=\bfseries\sffamily,
  title={Intuition},
  attach boxed title to top left={yshift=-2mm, xshift=4mm},
  boxed title style={colback=accentpurple!80, rounded corners},
  arc=4pt, boxrule=0.8pt,
  left=6pt, right=6pt, top=8pt, bottom=6pt,
  #1
}

% ============================================================
%  PAYOFF MATRIX HELPER
%  Usage: inside tabular, use \payoff{a}{b} for (a,b)
% ============================================================
\newcommand{\payoff}[2]{$(#1,\;#2)$}
\newcommand{\nash}[2]{\underline{$#1$},\underline{$#2$}}  % underline NE payoffs
\newcommand{\Nset}{\mathbb{N}}
\newcommand{\Rset}{\mathbb{R}}
\newcommand{\Sset}{\mathcal{S}}
\newcommand{\Ui}{u_i}

% ============================================================
%  SECTION FORMATTING
% ============================================================
\titleformat{\section}
  {\large\bfseries\sffamily\color{defblue}}
  {\thesection}{1em}{}[\titlerule]

\titleformat{\subsection}
  {\normalsize\bfseries\sffamily\color{thmgreen!80!black}}
  {\thesubsection}{1em}{}

\titleformat{\subsubsection}
  {\normalsize\itshape\sffamily\color{exorange!80!black}}
  {\thesubsubsection}{1em}{}

% ============================================================
%  HEADER / FOOTER
% ============================================================
\pagestyle{fancy}
\fancyhf{}
\fancyhead[L]{\sffamily\small\color{defblue} MTH636 -- Game Theory}
\fancyhead[R]{\sffamily\small\color{defblue} \leftmark}
\fancyfoot[C]{\sffamily\small\thepage}
\renewcommand{\headrulewidth}{0.4pt}

% ============================================================
%  HYPERREF SETUP
% ============================================================
\hypersetup{
  colorlinks=true,
  linkcolor=defblue,
  urlcolor=thmgreen,
  citecolor=accentpurple,
  bookmarksnumbered=true,
  pdftitle={MTH636 Game Theory Notes},
  pdfauthor={MTH636 Student Notes}
}

% ============================================================
%  MISC HELPERS
% ============================================================
\newcommand{\BR}{\text{BR}}        % Best Response
\newcommand{\NE}{\text{NE}}        % Nash Equilibrium
\newcommand{\SPNE}{\text{SPNE}}    % Subgame Perfect NE
\newcommand{\Dom}{\text{dom}}      % Dominance
\newcommand{\mixs}[1]{\sigma_{#1}} % Mixed strategy
\newcommand{\pures}[1]{s_{#1}}     % Pure strategy


% ============================================================
\begin{document}

% Title
\begin{center}
  {\LARGE\bfseries\sffamily\color{defblue} MTH636 -- Game Theory}\\[6pt]
  {\Large\sffamily Lecture 29: Strategic Form $\leftrightarrow$ Extensive Form \& Incomplete Information Games}\\[4pt]
\end{center}
\vspace{0.5cm}
\hrule height 1.5pt
\vspace{0.3cm}

% ============================================================
\section{From Strategic Form to Extensive Form}

\subsection{Motivation}

Suppose you have a strategic form game $G$.

\begin{itemize}
  \item \textbf{Question 1:} Can you get an equivalent extensive form game?
  \[
    G \;\longrightarrow\; \Gamma \;\longrightarrow\; G \quad (\text{equivalence})
  \]
  \item \textbf{Question 2:} If it exists, whether such an extensive form is unique or not?
\end{itemize}

\subsection{Answer to Question 1}

\noindent\textbf{Answer to Question 1:} Yes. Given
\[
  G = \bigl\langle N,\, (S_i)_{i \in N},\, (u_i)_{i \in N} \bigr\rangle,
\]
assume that $S_i = \{s_i^1, \ldots, s_i^{m_i}\}$ for each $i$.

\medskip
We construct an extensive form game $\Gamma$ as follows. The tree has the structure:
\begin{itemize}
  \item The root belongs to Player I. She has $m_1$ actions: $s_1^1, s_1^2, \ldots, s_1^{m_1}$.
  \item Each of Player I's $m_1$ children belongs to Player II (all in a single information set -- Player II cannot observe Player I's move). Player II has $m_2$ actions: $s_2^1, s_2^2, \ldots, s_2^{m_2}$.
  \item And so on for Player III, $\ldots$, Player $n$, each player having a single information set (so they cannot observe any previous moves).
  \item Every leaf node corresponds to a unique strategy profile in the strategic form game. The payoffs at that leaf are the utilities would $u_i$ of the corresponding profile.
\end{itemize}

\begin{notebox}
Every leaf node corresponds to a unique strategy profile in the strategic form game.
\end{notebox}

\medskip
It is easy to see that the corresponding strategic form game of the above extensive form game is $G$ only.

\subsection{Answer to Question 2 (Uniqueness)}

If you change the ordering of the players in the extensive form game, you will get another equivalent form.

\subsubsection*{Example}

Consider the following $2 \times 4$ strategic form game between Player 1 (rows $L$, $R$) and Player 2 (columns $rt$, $rb$, $\ell t$, $\ell b$):

\begin{center}
\begin{tabular}{c|cccc}
  1 $\backslash$ 2 & $rt$ & $rb$ & $\ell t$ & $\ell b$ \\\hline
  $L$ & $(2,-1)$ & $(2,-1)$ & $(2,0)$ & $(2,0)$ \\
  $R$ & $(2,1)$  & $(1,0)$  & $(2,1)$ & $(1,0)$ \\
\end{tabular}
\end{center}

This strategic form game can be represented as an extensive form game in (at least) two ways:

\paragraph{Extensive form with Player 1 moving first:}
\begin{itemize}
  \item Player I moves first: $L$ or $R$.
  \item After $L$: Player II is in an information set and chooses $r$ or $\ell$ (each leading to further choices).
  \item After $R$: Player II is in an information set and chooses $r$ or $\ell$, etc.
  \item Player II's strategy space consists of combinations $\{rt, rb, \ell t, \ell b\}$.
\end{itemize}

The resulting extensive form tree (Player I at root, Player II at second level in a single information set):
\begin{itemize}
  \item Root: Player I chooses $L$ or $R$.
  \item Left branch ($L$): Player II (single info set) chooses $r$ or $\ell$.
    \begin{itemize}
      \item[$r$:] leads to leaf with payoffs $(2,-1)$.
      \item[$\ell$:] leads to leaf with payoffs $(2,0)$.
    \end{itemize}
  \item Right branch ($R$): Player II (same info set? or separate) chooses $t$ or $b$ (or $r$ and $\ell$ leading to $t/b$).
\end{itemize}

\paragraph{Alternative extensive form (Player II depicted first at the single-info-set level):}

A cleaner representation places Player I at the root choosing $L$ or $R$, then Player II at two separate nodes (each a subgame given I's move), both enclosed in a single information set (since Player II does not observe Player I's choice). Terminal payoffs match the matrix entries.

\begin{center}
\begin{tabular}{c}
\textit{Extensive form tree (Player I first):}\\[6pt]
Player I: $L$ or $R$ \\
$\quad\Downarrow$ \\
Player II (information set): $n$ (right) or $\ell$ (left) \\
$\quad\Downarrow$ \\
Leaves: $(2,-1)$, $(2,0)$ after $L$; $(2,1)$, $(1,0)$ after $R$ with $r,\ell$ respectively.
\end{tabular}
\end{center}

\begin{notebox}
The extensive form representation of a strategic form game is \textbf{not unique}. Different orderings of the players yield different (but equivalent) extensive form games.
\end{notebox}

% ============================================================
\section{Incomplete Information Games}

\subsection{Motivation}

This model is due to \textbf{Harsanyi (1967)}.

\medskip
\noindent\textbf{Motivating Example.} Fix 13 cards. We can calculate the probabilities of the cards the other players are having:
\[
  \binom{52}{13}\binom{39}{13}\binom{26}{13}\binom{13}{13}
  \quad\text{or equivalently}\quad
  \binom{39}{13}\binom{26}{13}.
\]
Since each player's hand is drawn from the remaining cards, the type of each player (which 13 cards they hold) is private information -- other players only know the distribution.

\subsection{Definition}

\begin{defbox}{Incomplete Information Game}
An incomplete information game is a tuple:
\[
  \bigl\langle N,\, (T_i)_{i \in N},\, p,\, S,\, (s_A)_{A \in S \times T_i} \bigr\rangle
\]
where:
\begin{itemize}
  \item $N$: a finite set of players.
  \item $T_i$: a finite set of \textbf{types} for player $i$. The set of type vectors (type profiles) is denoted by $T = \bigtimes_{i \in N} T_i$.\\
    (In the card example, $|T_i| = \binom{52}{13}$; and $|T| = \left[\binom{52}{13}\right]^4$.)
  \item $p$: a probability distribution over $T$. We assume that
    \[
      p(t_i) \;=\; \sum_{t_{-i} \in T_{-i}} p(t_i, t_{-i}) \;>\; 0
      \qquad \forall\, i \in N,\;\; \forall\, t_i \in T_i.
    \]
    (Every type of every player is possible.)
  \item $S$: a set of \textbf{state games}. Every state game $s \in S$ is itself a vector:
    \[
      s = \bigl\langle N,\, (A_i(t_i))_{i \in N, t_i \in T_i},\, (u_i)_{i \in N} \bigr\rangle
    \]
    where $A_i$ is a non-empty set of actions available to player $i$, and $u_i : X \times A_i \to \mathbb{R}$ (where the product is taken over $i \in N$) is the utility function of player $i$.
  \item $s_A$: an element in $S$ -- the state game corresponding to a type profile.
\end{itemize}
\end{defbox}

\begin{notebox}
It is important to note that player $i$'s action depends on \textbf{his own type only}. Specifically, $s_A$ is an element in $S$ where:
\[
  s_A = \bigl\langle N,\, (A_i(t_i))_{i \in N,\, t_i \in T_i},\, (u_i(t))_{i \in N} \bigr\rangle.
\]
Player $i$'s action is chosen from $A_i(t_i)$, which depends on $t_i$ (player $i$'s own type) only.
\end{notebox}

\end{document}
